\input{estilo}
\begin{document}
\begin{figura}\begin{tikzpicture}[font=\small, piso/.style={draw=tinta, fill=white}]
  % la figura 4.7 del Saltzer
  \foreach \x/\t in {0/Client, 7.4/Service} {
    \draw[marco, rounded corners=8pt] (\x,0) rectangle ++(4.6,-5.6);
    \node[above] at (\x+2.3,0) {\t};
    \draw[piso] (\x+0.35,-0.35) rectangle ++(3.9,-2.3);
    \draw[piso] (\x+0.35,-2.95) rectangle ++(3.9,-2.3);
    \node at (\x+2.0,-4.0) {Stub};
  }
  \node[cod, anchor=north west] at (0.5,-0.55) {GET\_TIME()};
  \node[anchor=north west] at (1.25,-1.3) {10:42:18 PM EST};
  \node[cod, anchor=north west, align=left] at (7.9,-0.55) {GET\_TIME()};
  \node at (9.6,-1.25) {…};
  \node[anchor=north west] at (7.9,-1.5) {\textbf{return} \textit{time}};
  % el pedido baja, cruza y sube
  \draw[msg] (1.0,-1.0) -- (1.0,-4.7);
  \draw[msg] (1.6,-4.95) -- node[cod, above] {REQUEST} (10.4,-4.95);
  \draw[msg] (11.4,-4.7) -- (11.4,-1.8);
  % la respuesta vuelve
  \draw[msg] (9.0,-2.15) -- (9.0,-3.6);
  \draw[msg] (8.9,-3.8) -- node[cod, above] {RESPONSE} (2.5,-3.8);
  \draw[msg] (2.4,-3.6) -- (2.4,-2.2);
\end{tikzpicture}\end{figura}
\end{document}
